% ch11-universe — deals from companies we cannot see: (C6), (C2), the outside
% option, the external ground component, and honest accounting of unseen mass

\section{Two populations, one deal stream}\label{ch11:sec:populations}

Every model developed so far has quietly assumed that we know who could
receive the next deal. This chapter is about the two ways that assumption
fails. First, deals arrive from firms that are simply not in the tracked
universe: the deal row exists, carries a date and a sector, and is flagged
\code{company\_universe = 'Out of Universe'}, but the firm has no
companies-table row, no covariates, and no history --- defect (C6) of
\cref{ch:data}. Second, deals arrive from firms that will be in the
universe tomorrow but are invisible today: a firm enters observation
\emph{by raising}, so the tracked population is selected on the very event
we model --- defect (C2). Neither failure is small. In the synthetic
rehearsal, roughly $27\%$ of all deals went to firms outside the current
risk set, and the share sat at $23$--$25\%$ in held-out test windows
(\cref{app:results}). A mark factor that ignores this mass overstates
every tracked firm's probability by a factor of roughly $1.3$--$1.4$; a
ground model that ignores it cannot forecast sector totals honestly. The
remedy of this chapter is structural, not cosmetic: an explicit outside
option in the mark factor, an external component in the ground process, a
model-free audit of the unseen mass, and an entry model for the population
side of (C2).

We first fix what is and is not observed, in the notation of
\cref{def:ch02:market}.

\begin{definition}[Tracked and untracked populations; observed stream]\label{def:ch11:universes}
The firm universe splits as $\firms \cup \firms^{\dagger}$: the
\emph{tracked} universe $\firms$ (a companies-table row, covariates
$\tilde\xcov_i$, an event history $N_i$) and the \emph{untracked}
remainder $\firms^{\dagger}$ (no row, no covariates, no attributable
history). The observed deal stream is
\begin{equation}\label{eq:ch11:observed}
\big\{(t_m,\, s_m,\, \iota_m)\big\}_{m\ge 1},
\qquad
\iota_m \;=\;
\begin{cases}
i_m, & i_m \in \firms,\\[2pt]
\text{`Out of Universe'}, & i_m \in \firms^{\dagger},
\end{cases}
\end{equation}
together with deal marks $\dmark_m$: the time and sector of every deal are
recorded, but firm identity is recorded only on $\firms$. Operationally, a
deal is \emph{untracked} if its universe flag says so \emph{or} if its
firm makes its first-ever appearance in the data at $t_m$ (no prior
attributable history existed at $t_m^-$).
\end{definition}

\begin{assumption}[Complete sector stream]\label{asm:ch11:complete}
Every deal row carries a deal time and a sector label; no deal is dropped
from the stream for being untracked. Sector labels are correct up to the
small mislabel noise documented in \cref{app:columns}, whose treatment
(hygiene rules; the group-level variant of
\cref{ch11:sec:outside}) is orthogonal to universe membership.
\end{assumption}

\Cref{asm:ch11:complete} is a statement about the vendor's delivery
format, and it is the load-bearing asymmetry of this chapter: sector-level
ground counts are \emph{complete} --- every deal has a sector --- while
firm-level attribution is complete only on $\firms$. The consequence
deserves to be stated as a result, because it silently decides which
layers of the architecture (C6) can hurt.

\begin{proposition}[Attribution asymmetry]\label{prop:ch11:asymmetry}
Under \cref{asm:ch11:complete}, with the exact factorization of
\cref{thm:ch03:groundmark}:
\begin{enumerate}[label=(\roman*)]
\item The ground factor --- the likelihood of times and sectors, with
sector intensities $\Lam_s(t)$ built from sector-level histories and
exogenous covariates --- is a functional of fully observed data. Every
ground-process estimand and estimator is invariant to (C6): the ground
layer is immune.
\item The mark factor is not. Let $\pi_{s,t} = \Prob(i_{\text{next}}
\notin \riskset_s(t) \st \text{deal in } s \text{ at } t, \Ftm)$ be the
probability that the next deal in $s$ goes outside the tracked risk set.
A mark factor fitted on tracked-attributed deals only estimates the
conditional law $\Prob(i \st \text{deal}, i \in \riskset_s(t))$, and for
every tracked firm $i$,
\begin{equation}\label{eq:ch11:inflation}
\Prob\big(i \st \mathrm{deal\ in\ } s,t \big)
 \;=\;
 (1 - \pi_{s,t})\;
 \Prob\big(i \st \mathrm{deal\ in\ } s,t,\; i \in \riskset_s(t)\big).
\end{equation}
Reading the restricted model's output as an absolute probability
overstates every tracked firm by the factor $(1-\pi_{s,t})^{-1}$.
\item Any firm-level intensity assembled from the two layers,
$\lam_i(t) = \Lam_{s(i)}(t)\,\Prob(i \st \mathrm{deal})$, inherits the
inflation of (ii) multiplicatively, as do all horizon probabilities
(target T3) built from it.
\end{enumerate}
\end{proposition}

\begin{proof}
(i) By \cref{thm:ch03:groundmark} the ground factor of the log-likelihood
is $\sum_m \log \Lam_{s_m}(t_m) - \sum_s \int_0^{\Tend} \Lam_s(u)\dd u$,
which depends on the data only through the pairs $(t_m, s_m)$ and the
exogenous covariate paths. By \cref{asm:ch11:complete} these are observed
for every deal, untracked or not; hence the ground likelihood, its
maximizer, and any functional of the fitted $\Lam_s$ are unchanged by the
loss of firm identities. (The caveat --- that $\Lam_s$'s history term must
itself be built from \emph{all} deals --- is exactly the histories rule of
\cref{ch11:sec:ground}; violate it and immunity is forfeited.)
(ii) The event $\{i\} \subseteq \{i \in \riskset_s(t)\}$ for
$i \in \riskset_s(t)$, so the chain rule gives
$\Prob(i) = \Prob(i \in \riskset)\,\Prob(i \st i \in \riskset)
= (1-\pi_{s,t}) \Prob(i \st i \in \riskset)$. A likelihood evaluated only
on deals with $i_m \in \riskset$ is, term by term, the likelihood of the
conditional law given $\riskset$-membership; its MLE targets that
conditional law, not the absolute one.
(iii) Multiply (ii) by $\Lam_s(t)$; the marking theorem for marked point
processes \citep[Ch.~6]{daley2003} identifies
$\Lam_s(t)\Prob(i\st\mathrm{deal})$ as firm $i$'s intensity, and the
factor $(1-\pi_{s,t})^{-1}$ carries through every monotone functional of
$\lam_i$, including $p_i(t,\horizon)$.
\end{proof}

At the rehearsal's test-window shares, $\pi \approx 0.23$--$0.25$, the
inflation in \cref{eq:ch11:inflation} is $30$--$33\%$ on every tracked
firm --- uniform enough to leave within-sector \emph{rankings} intact
(which is why \cref{ch:ranking} could defer the issue) but fatal to any
absolute probability, any cross-sector comparison in which $\pi_{s,t}$
differs by sector, and any sector total decomposed by firm.
\Cref{prop:ch11:asymmetry} is also one more argument for the two-layer
architecture of \cref{thm:ch03:groundmark}: the defect ledger of
\cref{tab:ch02:defects} assigns (C6) to the mark layer \emph{alone},
so the timing deliverable (T2) survives on sector-complete counts while
the identity deliverable (T1) is repaired where it is broken --- in the
mark factor, by the outside option of \cref{ch11:sec:outside}.

\section{Newcomers are not entrants}\label{ch11:sec:newcomer}

Before building the repair, one distinction must be carved into the
chapter, because every downstream interpretation depends on it.

\begin{warning}[Newcomer $\ne$ entrant]\label{warn:ch11:entrant}
The \emph{newcomer probability} $\pi_{s,t}$ is the probability that the
next deal's firm is not in the current tracked risk set --- equivalently,
that the deal's firm has no attributable prior appearance in the data.
This is an \emph{observable}: every deal either matches a tracked history
or it does not, so $\pi_{s,t}$ can be estimated, calibrated, and audited.
The \emph{entrant probability} --- the probability that a genuinely new
economic entity enters the market --- is \emph{not identified} from this
data without an entry model and auxiliary coverage information, because a
first appearance conflates two channels: economic entry (a firm's first
raise ever) and first-time \emph{tracking} of an old firm (vendor
coverage catching up). The two channels are an unidentified mixture: in
the rehearsal, untracked incumbents accounted for roughly $12\%$ of
out-of-universe deals --- a nontrivial slice of the $\sim 27\%$ newcomer
share --- and
separating the channels with oracle labels changed the likelihood by only
$+0.003$ nats (\cref{app:results}). The conflation therefore costs almost
nothing in fit and cannot be detected by fit; what it costs, permanently,
is every compositional claim: ``entry accelerated in 2021,'' ``new
company formation shifted to fintech,'' ``the entrant cohort of 2020 is
larger than 2019's'' are all unsupported unless the coverage channel is
modeled or measured. Policy: in code, documentation, and reporting, this
repository says \emph{newcomer}, never \emph{entrant}, for the observable;
\emph{entrant} is reserved for the entry-model constructs of
\cref{ch11:sec:entry}, always with their identification caveats attached.
\end{warning}

The warning has a converse worth stating: because the two channels are
indistinguishable in likelihood, a single outside option that absorbs
both loses essentially nothing as a \emph{predictive} device. The
architecture below therefore models the newcomer mass as one bucket and
defers the decomposition to the entry model and the coverage audit, which
are the only places it can honestly live.

\section{The mark factor with an outside option}\label{ch11:sec:outside}

\subsection{The extended risk set}\label{ch11:sec:extended}

The repair to the mark factor is a one-row change with exact semantics.
Extend each sector's risk set with an \emph{outside option} $\lozenge_s$:
\begin{equation}\label{eq:ch11:extendedset}
\riskset^{+}_s(t) \;=\; \riskset_s(t) \cup \{\lozenge_s\},
\end{equation}
and give the outside option a score built from \emph{context} covariates
only --- nothing firm-level exists on that branch by construction:
\begin{equation}\label{eq:ch11:outsidescore}
\score_{\lozenge_s}(t) \;=\; v\T g_{s,t},
\qquad
g_{s,t} = \big[\text{trailing newcomer share},\ \text{sector deal volume},\
\text{macro as-ofs},\ \text{vintage/era effects},\ 1\big],
\end{equation}
while tracked firms keep their ranker scores $\score_i(t) =
\bcoef\T\zfeat_i(t)$ from \cref{ch:ranking}. The mark distribution is the
softmax over the extended set,
\begin{equation}\label{eq:ch11:extendedsoftmax}
\Prob\big(\dmark = a \st \text{deal in } s,t,\ \Ftm\big)
 \;=\;
 \frac{e^{\score_a(t)}}
      {e^{\score_{\lozenge_s}(t)} + \sum_{j \in \riskset_s(t)} e^{\score_j(t)}},
\qquad a \in \riskset^{+}_s(t),
\end{equation}
so that $\pi_{s,t} = \Prob(\lozenge_s \st \cdot)$ is produced by the same
softmax that produces firm probabilities. The next proposition says this
is not a heuristic patch but the \emph{correct} conditional model of the
observable mark when the next deal may go outside $\firms$.

\begin{proposition}[The collapsed mark model is exact]\label{prop:ch11:collapse}
Let $P^{*}(\cdot \st s,t,\Ftm)$ be the true next-deal firm distribution
on $\firms \cup \firms^{\dagger}$, and define the collapsed mark
$\tilde\dmark = i$ if $i \in \riskset_s(t)$ and $\tilde\dmark =
\lozenge_s$ otherwise. Then:
\begin{enumerate}[label=(\roman*)]
\item $\tilde\dmark$ is fully observed for \emph{every} deal, including
untracked ones; the extended-softmax likelihood built from
\cref{eq:ch11:extendedsoftmax} is the exact mark-factor likelihood of the
collapsed marked process --- no approximation is introduced by the
collapse, and no deal is dropped.
\item The family \cref{eq:ch11:extendedsoftmax} is well specified for the
collapsed law: for any categorical distribution with full support on
$\riskset^{+}_s(t)$ there exist scores realizing it (logits equal
log-probabilities up to an additive constant), and the parameterization
adds only the outside score \cref{eq:ch11:outsidescore}; identification
is up to the usual location shift, with the outside option playing the
role of the outside good's alternative-specific constant
\citep{mcfadden1974,train2009}.
\item Estimation is unchanged: the log-likelihood is the standard
conditional-choice form of \cref{ch:ranking} with one extra row per
choice set, and its MLE targets the KL projection of the true collapsed
law onto the family. In particular the tracked-firm probabilities it
outputs are \emph{absolute} --- the inflation of
\cref{eq:ch11:inflation} is removed by construction, since
$\Prob(i \st \cdot) = (1-\pi_{s,t}) \cdot
\mathrm{softmax}_{\riskset_s(t)}(\score)_i$ holds identically in the
extended model.
\end{enumerate}
\end{proposition}

\begin{proof}
(i) The collapse map $a \mapsto \tilde\dmark$ is a deterministic,
measurable coarsening of the mark space. For a tracked deal we observe
$i_m$ and hence $\tilde\dmark_m$; for an untracked deal we observe
precisely the event $\{i_m \notin \riskset_{s_m}(t_m)\}$, which \emph{is}
$\tilde\dmark_m = \lozenge_{s_m}$. The likelihood of a coarsened mark is
the pushforward of $P^*$ under the collapse, and
\cref{eq:ch11:extendedsoftmax} is a density for exactly that pushforward;
nothing unobservable enters. What is lost is only resolution
\emph{within} $\firms^{\dagger}$ and among out-of-pool firms --- which
the data does not carry anyway.
(ii) Given target probabilities $(p_a)_{a \in \riskset^+}$ with
$p_a > 0$, set $\score_a = \log p_a + c$: the softmax returns $(p_a)$
for any $c$, and conversely the softmax determines scores up to $c$
--- the standard bijection between positive categorical laws and score
differences. Restricting $\score_{\lozenge}$ to the span of $g_{s,t}$
makes the family parametric in the usual sense; whether the \emph{true}
$\pi_{s,t}$ lies in that span is an empirical question answered by the
calibration protocol below, not by the algebra.
(iii) Adding one alternative to a conditional-logit choice set changes
neither the objective's form, nor its concavity in the linear-in-features
parameterization, nor the softmax gradient; misspecified-model MLE
converges to the KL projection under the usual regularity. The displayed
identity is the chain rule applied inside the softmax:
conditioning \cref{eq:ch11:extendedsoftmax} on $\{a \ne \lozenge_s\}$
cancels the outside term from the denominator.
\end{proof}

Two estimation remarks. First, at newcomer shares of $23$--$27\%$ the
outside option is not a rare event, so the intercept bias of rare-event
logistic estimation \citep{kingzeng2001} is minor; if per-sector shares
in thin sectors fall to a few percent, the Firth correction
\citep{firth1993} costs one line and removes the concern. Second, the
context features in \cref{eq:ch11:outsidescore} are exactly the
covariates that made the newcomer share \emph{calibratable over time} in
the rehearsal: with the enriched context specification, quarterly
newcomer-share forecasts calibrated to within $\pm 2$pp on test
quarters, the standing acceptance target (\cref{app:results}). The trailing newcomer share
must be computed on an expanding or trailing window strictly before $t$
(\cref{prot:ch02:pit}), and the vintage/era terms exist to absorb the
coverage drift of \cref{ch11:sec:drift} rather than let it contaminate
the macro loadings.

\subsection{The architecture is a two-stage classifier, recognized}\label{ch11:sec:twostage}

The committed factorization --- ground process for sector timing, mark
factor for firm-given-sector with an outside option --- should be
recognized as an instance of a familiar design pattern: \emph{predict the
group first, then the member within the group, with an outlier class}.

\begin{proposition}[Two-stage equivalence]\label{prop:ch11:twostage}
Given a deal at $t$, let $\Lam(t) = \sum_{s} \Lam_s(t)$ be the
superposed ground intensity. Then for any tracked firm $i$ with sector
$s(i)$,
\begin{equation}\label{eq:ch11:chain}
\Prob\big(i \st \mathrm{deal\ at\ } t, \Ftm\big)
 \;=\;
 \underbrace{\frac{\Lam_{s(i)}(t)}{\Lam(t)}}_{\mathrm{stage\ 1:\ sector}}
 \;\times\;
 \underbrace{\frac{e^{\score_i(t)}}
   {e^{\score_{\lozenge_{s(i)}}(t)} + \sum_{j\in\riskset_{s(i)}(t)} e^{\score_j(t)}}}_{\mathrm{stage\ 2:\ firm\ or\ outlier}},
\end{equation}
and the residual mass $\sum_s \frac{\Lam_s(t)}{\Lam(t)}\pi_{s,t}$ is the
probability the deal goes outside the tracked universe altogether. Thus
the two-layer model \emph{is} the two-stage sector$\to$firm classifier
with a per-sector outlier bucket; conversely, any such two-stage
classifier with strictly positive stage probabilities defines a valid
mark distribution on $\bigcup_s \riskset^{+}_s(t)$.
\end{proposition}

\begin{proof}
By the marking theorem for the superposition of sector streams
\citep[Ch.~6]{daley2003}, conditional on a deal at $t$ the sector is
distributed as $\Lam_s(t)/\Lam(t)$; multiplying by the conditional firm
law \cref{eq:ch11:extendedsoftmax} is the chain rule. The converse
direction is immediate: nonnegative stage probabilities that sum to one
compose, by the same chain rule, into a categorical law on the disjoint
union of extended risk sets.
\end{proof}

The equivalence earns its display for one practical reason: it makes the
\emph{group-level variant} a first-class alternative rather than an
architectural upheaval. Sector labels carry mislabel noise
(\cref{asm:ch11:complete}; \cref{app:columns}), and at the finest
granularity a mislabeled deal both pollutes one sector's ground counts
and lands in the wrong choice set. Collapsing sectors into industry
groups is another measurable coarsening of the mark --- the same quotient
construction as \cref{prop:ch11:collapse} --- so stage 1 may equally
predict the industry \emph{group} and stage 2 rank firms within the
group's (larger) risk set, with one outlier bucket per group. Mislabels
that stay within a group then cancel entirely. The trade is variance for
bias: coarser groups mean larger softmax denominators and weaker
sector-specific context in \cref{eq:ch11:outsidescore}. The decision
between granularities belongs to the goodness-of-fit machinery of
\cref{ch:gof} applied per level; the mathematics is level-agnostic,
provided --- and this is a hard constraint --- ground and mark layers use
the \emph{same} partition, or \cref{eq:ch11:chain} silently stops being
a probability distribution.

\subsection{Calibration of the outside option is a deliverable}\label{ch11:sec:calibration}

With the outside option in place, the newcomer share stops being a
nuisance and becomes a first-class model output: $\pi_{s,t}$ is a
predicted probability, realized outcomes are observed for every deal, and
the pair supports exactly the calibration discipline of \cref{ch:gof}.

\begin{protocol}[Reliability of the outside option]\label{prot:ch11:reliability}
On every evaluation window: (a) per quarter and per sector, compare the
average predicted $\pi_{s,t}$ over deals against the realized newcomer
share, reporting the mean absolute error (rehearsal reference: $3$--$4$pp
with context covariates, $6$--$9$pp for a constant;
\cref{app:results}); (b) pool deals into bins by predicted $\pi$ and
report a reliability diagram with binomial bands; (c) report the same by
era, so that coverage drift (\cref{ch11:sec:drift}) is visible as an
era-structured miscalibration rather than averaged away. Failures of (a)
with stable context covariates are escalated to the coverage audit
(\cref{prot:ch11:coverage}) \emph{before} any model change.
\end{protocol}

\begin{decision}{D11.1 --- outside option in the next ranker revision}
The next revision of the mark-factor ranker extends every choice set with
the outside-option row \cref{eq:ch11:extendedset}--\cref{eq:ch11:outsidescore}:
context features on the $\lozenge_s$ row only, firm features on tracked
rows only, one shared softmax. This extends
\modrepo{sector\_ranker.py}'s risk-set construction and mirrors the
synthetic loader's \code{build\_choice\_sets(newcomer\_context=True)}
path, with \code{experiments/exp28\_newcomer\_designs.py} (designs
N1--N4) as prior art. Rationale: \cref{prop:ch11:asymmetry} (the
restricted model's probabilities are inflated by $(1-\pi)^{-1}$),
\cref{prop:ch11:collapse} (the extension is exact and costs nothing in
estimation), and the measured calibration gain of context covariates.
Reversal trigger: if on real data the realized newcomer share falls below
$2\%$ in every sector and era \emph{and}
\cref{prot:ch11:reliability} shows a constant share calibrates within
1pp, the outside option may be frozen to an intercept-only score ---
removed it is not, since \cref{eq:ch11:inflation} holds at any $\pi>0$.
\end{decision}

\section{The ground process: the external component and the histories rule}\label{ch11:sec:ground}

\subsection{Decomposing the immigrant intensity}\label{ch11:sec:mudecomp}

On the ground layer, (C6) and (C2) appear not as missing labels but as a
component of intensity that no tracked-firm bookkeeping can generate.
Write the sector ground intensity in Hawkes form
(\cref{ch:hawkes}) and split the baseline:
\begin{equation}\label{eq:ch11:mudecomp}
\Lam_s(t)
 \;=\;
 \underbrace{\mub_s^{\mathrm{tr}}(t) + \mub_s^{\mathrm{ext}}(t)}_{\mub_s(t)}
 \;+\; \sum_{m:\, t_m < t} \kernel_{s s_m}(t - t_m),
\end{equation}
where $\mub_s^{\mathrm{tr}}$ carries the part of the immigrant intensity
attributable to the tracked population's state and
$\mub_s^{\mathrm{ext}}$ is the \emph{external component}: deal flow from
firms we cannot see, driven by the context covariates $g_{s,t}$ of
\cref{eq:ch11:outsidescore} (and, once \cref{ch11:sec:entry} exists, by
cohort-based newcomer supply). The decomposition is not identified by the
ground layer alone --- both components are baselines --- but it is pinned
by the mark layer through the marking theorem:

\begin{corollary}[Newcomer substream intensity]\label{cor:ch11:split}
If the outside-option probability $\pi_{s,t}$ of
\cref{eq:ch11:extendedsoftmax} is predictable, then the newcomer deals of
sector $s$ form a point process with conditional intensity
$\pi_{s,t}\,\Lam_s(t)$, and the tracked-attributed deals one with
intensity $(1-\pi_{s,t})\,\Lam_s(t)$. The sector total decomposes
exactly as
\begin{equation}\label{eq:ch11:newintensity}
\Lam_s(t) \;=\; \underbrace{(1-\pi_{s,t})\Lam_s(t)}_{\text{tracked}}
\;+\; \underbrace{\pi_{s,t}\Lam_s(t)}_{\text{newcomer}},
\end{equation}
which is the position-dependent thinning of the ground process by the
mark factor \citep[Ch.~5]{daley2003}.
\end{corollary}

\Cref{eq:ch11:newintensity} is the honest version of a sector forecast:
a predicted quarterly total \emph{with its newcomer fraction attached},
rather than a tracked-only total silently presented as the market. It is
also the Hawkes--Oakes immigrant/offspring split
\citep{hawkesoakes1974} read at the mark level: the newcomer substream is
the immigrant channel of the branching representation, arriving from
outside the tracked system's genealogy. One rehearsal finding tempers the
interpretation: when the two substreams were fitted separately, the
\emph{newcomer} stream carried the larger apparent excitation radius
($0.30$ against $0.08$--$0.11$ for repeats; \cref{app:results}) --- entry
waves are where latent common factors live, so apparent ``excitation'' on
the newcomer stream must be read as clustering of unknown origin, in line
with the frailty-inflation verdict of \cref{ch:hawkes}
(\cref{app:results}: fitted branching ratios are upper bounds).

\subsection{The histories rule, and the price of breaking it}\label{ch11:sec:histories}

The excitation term in \cref{eq:ch11:mudecomp} raises the operational
question with the sharpest teeth in this chapter: \emph{which deals enter
the history sum?} Untracked deals are real events; they move the sector's
state and excite future deal flow exactly as tracked deals do. The rule
is therefore:

\begin{decision}{D11.2 --- histories include all deals}
Every history functional in every fitter --- lagged sector counts in the
weekly GLM, exponential-kernel sums in the event-time fitters, trailing
volume covariates, cooldown clocks at sector level --- is computed from
the \emph{complete} deal stream, untracked deals included. Only
\emph{firm-level attribution} (risk-set membership, firm histories
$N_i$, firm covariates) is restricted to $\firms$. Rationale:
\cref{prop:ch11:thinning} below --- dropping untracked deals from
histories understates excitation by a factor of about the tracked share
and distorts the decay estimate --- and \cref{prop:ch11:asymmetry}(i),
whose immunity claim is conditional on exactly this rule. Reversal
trigger: none; this is an accounting identity about what an event is, not
a tunable choice.
\end{decision}

The bias when the rule is broken can be derived exactly in the linear
case. Model the loss of untracked deals as independent thinning: each
event is retained (tracked) with probability $q = 1-p$ independently of
everything else --- optimistic, since real coverage is not independent of
firm type, but sufficient to sign and size the effect.

\begin{proposition}[Thinning bias in the branching ratio]\label{prop:ch11:thinning}
Let $N$ be a stationary linear Hawkes process with immigrant rate
$\mub$, kernel $\kernel$ with branching ratio
$\branch = \int_0^\infty \kernel < 1$, and let $N^{q}$ be its
independent thinning with retention $q = 1-p$
\citep[Ch.~11]{daley2003}. Then:
\begin{enumerate}[label=(\roman*)]
\item (Rates) $N^q$ has mean rate
$\bar\lam^{q} = q\,\mub/(1-\branch)$: thinning scales the mean by $q$
exactly.
\item (First generation) In the Hawkes--Oakes cluster representation
\citep{hawkesoakes1974}, the expected number of \emph{retained} direct
offspring of a retained event is $q\branch = (1-p)\branch$. Any
estimator that identifies branching by attributing observed events to
observed parents --- the declustering/EM estimand
\citep{veenschoenberg2008} --- therefore targets first-generation mass
scaled by $(1-p)$.
\item (All generations) The expected number of retained descendants (all
generations) of a retained event is $q\branch/(1-\branch)$; matching the
total-progeny identity $\branch_{\mathrm{app}}/(1-\branch_{\mathrm{app}})
= q\branch/(1-\branch)$ gives
\begin{equation}\label{eq:ch11:kappaapp}
\branch_{\mathrm{app}} \;=\; \frac{(1-p)\,\branch}{1 - p\branch}
 \;\in\; \big[(1-p)\branch,\ \branch\big),
\end{equation}
again $(1-p)\branch$ to first order in $\branch$. In both readings the
apparent branching ratio is understated, and the missing excitation mass
reappears as an inflated apparent baseline:
$\mub_{\mathrm{app}} = \bar\lam^{q}(1-\branch_{\mathrm{app}})
= q\mub/(1-p\branch) > q\mub$.
\end{enumerate}
\end{proposition}

\begin{proof}
(i) The stationary mean of a subcritical linear Hawkes solves
$\bar\lam = \mub + \branch\bar\lam$, giving $\mub/(1-\branch)$ (sum the
geometric cascade of generations, $\mub\sum_{k\ge0}\branch^k$).
Independent thinning multiplies the first moment measure by $q$
\citep[Ch.~11]{daley2003}.
(ii) In the cluster representation each event, independently of
everything else, produces direct offspring as an inhomogeneous Poisson
process with mean $\branch$. The retention marks are i.i.d.\
Bernoulli$(q)$, independent of the genealogy and of each other; by
Wald's identity the expected number of retained direct offspring is
$q\branch$, and conditioning on the parent's own retention changes
nothing by independence of the marks.
(iii) Generation-$k$ descendants of a fixed event have expected count
$\branch^k$; each is retained with probability $q$ independently, so the
expected retained progeny is $q\sum_{k\ge1}\branch^k =
q\branch/(1-\branch)$. Solving the displayed identity for
$\branch_{\mathrm{app}}$ gives \cref{eq:ch11:kappaapp}; monotonicity in
$q$ and the bounds follow by inspection ($q=1$ recovers $\branch$). The
baseline claim is (i) combined with
$\mub_{\mathrm{app}} = \bar\lam^q(1-\branch_{\mathrm{app}})$ and
$1-\branch_{\mathrm{app}} = (1-\branch)/(1-p\branch)$.
\end{proof}

Two remarks complete the picture. First, the \emph{decay} is distorted
alongside the scale: a retained grandchild whose intermediate parent was
dropped presents to the fitter as a direct child at the \emph{sum} of two
kernel lags, so the apparent kernel is contaminated by convolution powers
$\kernel^{*k}$ --- heavier-tailed than $\kernel$ --- and a fitted
exponential decay $\decay$ is dragged toward slower values. Since
\cref{ch:gof} already shows $\decay$ is the fragile parameter
(the $\branch$--$\decay$ ridge), history thinning attacks the model
exactly where it is weakest. Second, the repository's actual failure mode
is the \emph{mixed} case --- sector counts complete
(\cref{asm:ch11:complete}) but the history term built from tracked deals
only, as happens whenever a fitter joins deals to the companies table and
drops non-matches. Then the excitation covariate is a $q$-scaled, noisy
version of the true excitation state: an errors-in-variables regression
that attenuates the excitation coefficient and reroutes the missing mass
into the baseline and into any covariate correlated with activity --- the
same understatement, arrived at by a different door. The
covariates-dominate verdict of \cref{app:results} already leaves
excitation only decimal points of likelihood to live on;
\cref{prop:ch11:thinning} says that breaking the histories rule takes a
$(1-p)$-sized bite out of precisely those decimal points.

\section{Unseen mass: the model-free audit}\label{ch11:sec:unseen}

Everything so far trusts the model to measure the outside world through
$\pi_{s,t}$. Good practice demands an estimate of the unseen mass that
does not pass through the model at all, so that model and audit can
disagree informatively. The classical tool is exactly fitted to our
question: \emph{what is the probability that the next deal comes from a
firm we have not seen before?}

Fix a sector and a window containing $N$ deals, and regard the firm
identities as draws $X_1,\dots,X_N$ from an unknown distribution
$(p_i)$ over the (countable, partly invisible) firm population. Let
$N_r$ be the number of firms appearing exactly $r$ times in the window,
and let the \emph{missing mass} $M_0 = \sum_i p_i \ind{i \notin
\{X_1,\dots,X_N\}}$ be the probability mass of firms not seen at all.
The Good--Turing estimator \citep{good1953} is
\begin{equation}\label{eq:ch11:gt}
\widehat M_0 \;=\; \frac{N_1}{N},
\end{equation}
the fraction of deals going to firms seen exactly once.

\begin{proposition}[Good--Turing is nearly unbiased]\label{prop:ch11:gt}
If $X_1,\dots,X_N$ are i.i.d.\ (exchangeability suffices), then
\[
\E\big[\widehat M_0\big] \;=\; \E\big[M_0^{(N-1)}\big]
\qquad\text{and}\qquad
0 \;\le\; \E\big[\widehat M_0\big] - \E\big[M_0^{(N)}\big]
 \;=\; \sum_i p_i^2 (1-p_i)^{N-1}
 \;\le\; \tfrac{1}{N}\big(1-\tfrac1N\big)^{N-1} \;\le\; \tfrac{1}{2N},
\]
where $M_0^{(n)}$ is the missing mass of a sample of size $n$: the
estimator is unbiased for the missing mass of a sample one draw smaller,
hence biased upward for $M_0^{(N)}$ by at most $O(1/N)$.
\end{proposition}

\begin{proof}
Leave one out. By exchangeability,
\[
\E\Big[\frac{N_1}{N}\Big]
 = \frac1N \sum_{n=1}^{N} \Prob\big(X_n \text{ is a singleton in } X_{1:N}\big)
 = \Prob\big(X_N \notin \{X_1,\dots,X_{N-1}\}\big),
\]
since $X_N$ is a singleton in the full sample precisely when its firm is
absent from the other $N-1$ draws. Conditioning on the first $N-1$
draws, that probability is $\E[\sum_i p_i \ind{i \notin X_{1:N-1}}] =
\E[M_0^{(N-1)}]$, which proves the identity; the estimator answers
``will the \emph{next} deal come from an unseen firm'' almost by
definition. For the bias, $\E[M_0^{(n)}] = \sum_i p_i(1-p_i)^{n}$, so
\[
\E\big[\widehat M_0\big] - \E\big[M_0^{(N)}\big]
 = \sum_i p_i\big[(1-p_i)^{N-1} - (1-p_i)^{N}\big]
 = \sum_i p_i^2 (1-p_i)^{N-1}
 \le \max_{x\in[0,1]} x(1-x)^{N-1} \sum_i p_i,
\]
and the maximum is attained at $x = 1/N$ with value
$\frac1N(1-\frac1N)^{N-1} \le \frac1{2N}$ for $N\ge2$
(and $\to (eN)^{-1}$). When exchangeability fails the identity fails
with it: this is the assumption to watch, not the algebra.
\end{proof}

The proof makes the operative assumption explicit: exchangeability of
deals \emph{within the window}. On this data that is violated by
everything this document models --- excitation, macro dynamics, coverage
drift --- so the estimator is applied on \emph{rolling windows per
sector}, short enough that within-window dynamics are second-order,
long enough that $N_1$ is not noise; the smoothing practice of
\citet{galesampson1995} applies when small counts make the raw
frequencies-of-frequencies ragged. Two further caveats. The window
convention must match the functional: $N_1/N$ on a trailing window
estimates the window-relative unseen mass, which is the right comparator
for a model probability evaluated under the same convention ---
first-ever appearance versus first-in-window appearance are different
events, and the audit must fix one and compute both sides accordingly.
And GT estimates a \emph{probability mass}, not a population count. For
the count question --- how many fundable firms are we not seeing? --- the
Chao lower bound \citep{chao1984} on the number of distinct firms $S$
active in the window is
\begin{equation}\label{eq:ch11:chao}
\widehat S_{\mathrm{Chao}} \;=\; S_{\mathrm{obs}} + \frac{N_1^2}{2N_2}
 \;\le\; S \quad\text{(in expectation, as a lower bound)},
\end{equation}
valid as a floor precisely under the abundance heterogeneity that is
certain to hold here (firms differ wildly in deal propensity), which is
the honest direction: the invisible pool is \emph{at least} this large.

\begin{protocol}[Good--Turing versus the model]\label{prot:ch11:gtaudit}
Standing diagnostic, run with every evaluation campaign and reported
alongside \cref{prot:ch11:reliability}: per sector and rolling window,
(a) compute $\widehat M_0 = N_1/N$ and $\widehat S_{\mathrm{Chao}}$ from
the raw deal stream --- no model involved; (b) compute the model's
average predicted outside-option mass over the same deals under the same
window convention; (c) report the paired series and their gap. Agreement
within the binomial noise band is a nontrivial consistency check
(two estimators of one functional, one parametric and one not);
a drifting gap with stable context covariates is escalated to the
coverage audit of \cref{prot:ch11:coverage} before any refit.
\end{protocol}

\begin{decision}{D11.3 --- unseen-mass audit ships with the evaluation harness}
The GT/Chao audit of \cref{prot:ch11:gtaudit} is implemented once, in the
shared evaluation layer (planned \modrepo{eval/unseen\_mass.py}), and
run by default in every campaign --- not as an occasional notebook.
Rationale: it is the only estimate of the outside world that does not
share assumptions with the model it audits, and it is nearly free
(counting singletons). Reversal trigger: none anticipated; if window
choice proves contentious, the audit gains a sensitivity panel over
window lengths rather than being dropped.
\end{decision}

\section{Modeling entry: cohorts on the Lexis diagram}\label{ch11:sec:entry}

Defect (C2) has appeared so far as a missing-mass problem. Its population
face is different: the tracked universe is \emph{selected on having
raised}, so the data is silent about the reservoir of founded-but-not-yet-funded
firms from which tomorrow's newcomers arrive. An \emph{entrant-profile
module} models what can be modeled about that reservoir: the intensity of
\emph{first observed rounds}, organized by founding cohort.

The natural coordinate system is the Lexis diagram \citep{keiding1990}:
calendar time $t$ on one axis, firm age $a = t - c$ (time since founding,
cohort $c$) on the other; firms travel along unit-slope lines, and a
first observed round at age $a$ is a point on that line --- observed only
if it falls inside the data window (firms whose first raise postdates
June 2024 are invisible entirely: right truncation of entry, the
mirror image of (C1)). Because the number of not-yet-funded firms in a
cohort is unobserved, only the \emph{aggregate} first-round intensity is
identified, and that is what the module fits:
\begin{equation}\label{eq:ch11:cohort}
\lam^{\mathrm{entry}}_{s,c}(t)
 \;=\;
 \exp\!\big(\alpha_s + \phi_{\mathrm{era}(c)} + f_s(t - c)
 + \bcoef\T \xcov^{\mathrm{mac}}_{\mathrm{pub}}(t)\big),
\qquad
\mub_s^{\mathrm{ext}}(t) \;\supseteq\; \sum_{c \le t} \lam^{\mathrm{entry}}_{s,c}(t),
\end{equation}
with $f_s$ an age profile (entry hazards peak in the first years after
founding), $\phi$ era-of-founding effects (the dry-powder vintages), and
macro as-ofs on the calendar axis. Founding dates come from the
companies table --- available, with (C4)-grade noise, for firms once they
are tracked --- so the cohort structure is estimable from each firm's
\emph{first appearance}, the one moment at which entry-relevant features
are observed.

What the module buys is exactly two things. First, \emph{honest sector
totals}: with \cref{eq:ch11:cohort} feeding
$\mub_s^{\mathrm{ext}}$, the decomposition
\cref{eq:ch11:newintensity} becomes a forecast of incumbent
\emph{plus} newcomer intensity with the newcomer part driven by cohort
supply rather than extrapolated share --- the difference matters whenever
era effects move (a 2021-vintage cohort bulge produces newcomer flow for
years). Second, \emph{cohort-based forecasts of newcomer supply}:
questions like ``how much first-round flow should sector $s$ expect next
year given the founding activity of 2022--2024'' become answerable, with
the era effects carrying the dry-powder cycle.

What it cannot buy must be stated with equal clarity: \emph{firm-level
prediction for unseen firms}. A firm with no covariates admits no
firm-level score; no entry model changes that. The honest statement of
the module's output is a \emph{sector-level newcomer intensity} together
with a mark distribution over \emph{observable-at-entry} features ---
stage, geography, sector, and the deal-time profile measured at first
appearance --- learned from tracked firms at their first rounds. That
learned distribution is a \emph{selection} distribution (features
conditional on having entered), which supports predicting the profile of
the next entrant; it says nothing about the latent population of firms
that have not raised and may never (\cref{warn:ch11:entrant} applies in
full: the ``entry'' intensity \cref{eq:ch11:cohort} is estimated from
first \emph{appearances} and inherits the coverage channel unless the
audit of \cref{ch11:sec:drift} clears it). The rehearsal adds a
caution: newcomer utility was dominated by the alternative-specific
constant and the context/pool-size terms rather than by entrant-profile
quality (\cref{app:results}: ASC $6.86$ on regime A against
observable-profile weights an order of magnitude smaller) --- effort
belongs in $g_{s,t}$ and the cohort structure, not in rich profile
models.

\begin{decision}{D11.4 --- a planned \code{entrants/} module, scoped by identification}
The repository plans an \modrepo{entrants/} module implementing
\cref{eq:ch11:cohort}: cohort--age--era first-round intensities per
sector, an observable-at-entry mark distribution estimated at first
appearances, and the coupling of both into
$\mub_s^{\mathrm{ext}}$ and the outside-option features $g_{s,t}$. The
module's API exposes sector-level newcomer intensity and entrant-profile
distributions only --- it deliberately has no per-firm scoring surface,
so the identification limit is enforced by the interface rather than by
discipline. Rationale: the honest-totals and cohort-supply gains above;
the interface constraint is \cref{warn:ch11:entrant} made structural.
Reversal trigger: acquisition of auxiliary coverage data (a second
vendor, or vendor backfill metadata) that identifies the
entrant-versus-untracked split would justify promoting the module to a
genuine entry model with compositional claims.
\end{decision}

\section{Two-sided universe drift}\label{ch11:sec:drift}

One final structural fact cuts across everything above: the universe
boundary itself moves, from both sides. Firms cross into $\firms$ as
PitchBook's coverage expands, and coverage has expanded materially over
2011--2024 --- tracking improves, backfill happens, whole segments
(regions, stages) arrive in waves. The untracked share is therefore
time-varying \emph{for reasons that are not market dynamics}: a declining
\code{off\_universe\_share} is what improving coverage looks like, and it
is also exactly what a genuine shift toward incumbent-dominated deal flow
would look like. Every quantity in this chapter --- $\pi_{s,t}$, the GT
mass, the cohort intensities, the era effects of
\cref{eq:ch11:cohort} --- confounds the two explanations unless coverage
is audited first.

\begin{protocol}[Coverage audit by era]\label{prot:ch11:coverage}
The loader reports \code{off\_universe\_share} by sector and by era
(fixed calendar blocks) as a first-class output, on synthetic and real
data alike (\cref{ch:data} already requires the aggregate; this protocol
requires the era breakdown). Any trend in the series is treated as a
coverage artifact \emph{candidate} until examined: check for step
changes aligned with vendor events rather than market events, for
divergence between the untracked share and the GT newcomer mass
(\cref{prot:ch11:gtaudit}) --- which respond differently to backfill ---
and for era-localized failures of \cref{prot:ch11:reliability}.
\end{protocol}

\begin{decision}{D11.5 --- coverage audit precedes entry claims}
No claim about entry dynamics --- trends in newcomer share, cohort sizes,
era effects in \cref{eq:ch11:cohort} --- is reported, internally or
externally, until \cref{prot:ch11:coverage} has been run on the relevant
eras and the trend has survived it. Rationale: the drifting-share
tripwire is the one diagnostic that doubles as a data-quality monitor;
its most likely cause on this vendor's history is coverage, not the
market, and \cref{warn:ch11:entrant} already establishes that the data
cannot arbitrate by itself. Reversal trigger: none; this is an ordering
constraint on analysis, and it costs one audit per campaign.
\end{decision}

\section{Defect declarations and the ranking interface}\label{ch11:sec:defects}

Per \cref{tab:ch02:defects} and decision D2.1, this chapter \emph{owns}
(C2) and (C6), and its models relate to the ledger as follows.
\emph{(C6)}: the ground layer is robust under
\cref{asm:ch11:complete} and the histories rule
(\cref{prop:ch11:asymmetry}); the extended mark factor \emph{corrects}
the defect exactly at the collapsed-mark resolution
(\cref{prop:ch11:collapse}); fitters that break the histories rule are
vulnerable with signed bias --- excitation understated by a factor
$\approx(1-p)$, baseline inflated, decay dragged slow
(\cref{prop:ch11:thinning}). \emph{(C2)}: the outside option absorbs
its deal-stream face; the entry module models its population face at
sector level under the stated identification limits; compositional
claims remain vulnerable permanently absent auxiliary coverage data
(\cref{warn:ch11:entrant}), and the vulnerability is accepted and
fenced by D11.4/D11.5 rather than assumed away. \emph{(C1)}: robust ---
all windows and spells here inherit the censoring treatment of
\cref{ch:data,ch:censoring}; right truncation of entry is handled by
construction in \cref{eq:ch11:cohort}. \emph{(C3)}: not corrected here,
and the interaction matters: unlabeled dead firms \emph{deflate}
tracked-firm probabilities by padding the softmax denominator, while a
missing outside option \emph{inflates} them
(\cref{eq:ch11:inflation}); the two biases can partially cancel in
aggregate metrics, and this document forbids reading such cancellation
as health --- the two defects are diagnosed separately
(\cref{ch:ranking} for (C3), \cref{prot:ch11:reliability} here).
\emph{(C4)}: sector mislabels perturb \cref{asm:ch11:complete}
marginally and motivate the group-level variant of
\cref{ch11:sec:twostage}; date jitter is immaterial to shares at weekly
resolution. \emph{(C5)}: the outside-option score
\cref{eq:ch11:outsidescore} uses only context covariates, which are
never stale; the tracked rows inherit the treatment of \cref{ch:stale}
unchanged.

The interface with the ranking metrics of \cref{ch:ranking} must be
fixed explicitly, because the outside option changes what recall can
mean. Newcomer deals have no tracked target to recall: including them in
the denominator of recall@$K$ mechanically caps the metric at
$1-\pi$ and makes scores incomparable across eras with different
coverage. The coordinated protocol is therefore: recall@$K$ and MRR are
computed on tracked-attributed deals only, exactly as \cref{ch:ranking}
specifies; the newcomer mass is reported \emph{alongside} as the
calibrated $\pi_{s,t}$ output under \cref{prot:ch11:reliability}; and
every ranking report carries both numbers, so that a model cannot buy
apparent recall by silently wishing the outside world away, nor be
penalized for deals no tracked-universe model could have named.

\begin{codehooks}
\emph{Mark factor.} \modrepo{sector\_ranker.py} gains the extended risk
set \cref{eq:ch11:extendedset}: one outside-option row per choice set
carrying the context features \cref{eq:ch11:outsidescore}, mirroring
\code{synthetic/loaders.py}'s
\code{build\_choice\_sets(newcomer\_context=True)}; prior art and
reference fits in \code{experiments/exp28\_newcomer\_designs.py}
(N1 context score, N2 entrant moments, N3 two-step, N4
immigrant/offspring split).
\emph{Ground process.} The histories rule (D11.2) binds
\modrepo{sector\_stability.py}, the event-time fitters in
\modrepo{eventtime/}, and the interval-censored tools in
\modrepo{mbpp/}: history terms consume the complete deal stream; the
external baseline component of \cref{eq:ch11:mudecomp} enters as
covariates until \modrepo{entrants/} exists.
\emph{Audits.} \modrepo{eval/unseen\_mass.py} (to be written) implements
\cref{prot:ch11:gtaudit} (Good--Turing mass, Chao floor, model
comparison) and the reliability report of \cref{prot:ch11:reliability};
\modrepo{data/audits.py} extends its universe-share check to the by-era
breakdown of \cref{prot:ch11:coverage}.
\emph{Entry.} The planned \modrepo{entrants/} module implements the
cohort--age--era intensities \cref{eq:ch11:cohort} and the
observable-at-entry mark distribution, exposing sector-level outputs
only (D11.4).
\end{codehooks}
